\documentclass[a4paper]{article}
\usepackage[a4paper,margin=0mm]{geometry}
\usepackage{fontspec}
\usepackage{unicode-math}
\usepackage{tikz}
\usepackage[table]{xcolor}
\usepackage[absolute]{textpos}
\usepackage{array}

\setmainfont{Libertinus Serif}
\setsansfont{Libertinus Sans}
\setmathfont{Libertinus Math}

\definecolor{accent}{HTML}{F46530}
\definecolor{badge}{HTML}{FFBDA5}
\definecolor{ink}{HTML}{2A2E30}

\setlength{\TPHorizModule}{1mm}
\setlength{\TPVertModule}{1mm}
\setlength{\parindent}{0pt}
\pagestyle{empty}

\newcommand{\topicheader}[2]{%
\begin{textblock*}{194mm}(16mm,10mm)
\begin{tikzpicture}
\node[fill=accent,text=white,font=\sffamily\bfseries\fontsize{11}{11}\selectfont,
      minimum width=9mm,minimum height=9mm,inner sep=0pt] at (4.5mm,-4.5mm) {#1};
\node[font=\sffamily\bfseries\fontsize{15}{16}\selectfont,text=ink,anchor=west] at (12mm,-4.5mm) {#2};
\end{tikzpicture}
\end{textblock*}
\begin{textblock*}{194mm}(16mm,22.5mm)
\begin{tikzpicture}
\fill[accent] (0,0) rectangle (194mm,0.5mm);
\end{tikzpicture}
\end{textblock*}
}

\newcommand{\pagekicker}[1]{%
\begin{textblock*}{150mm}(16mm,25.5mm)
\begin{tikzpicture}
\node[fill=ink,text=white,font=\sffamily\bfseries\fontsize{8}{8}\selectfont,
      inner sep=2.2pt,rounded corners=1pt] {\MakeUppercase{#1}};
\end{tikzpicture}
\end{textblock*}
}

\newcommand{\pagefooter}[1]{%
\begin{textblock*}{194mm}(8mm,286.41mm)
  \begin{tikzpicture}
    \fill[accent] (0,0) rectangle (194mm,0.6mm);
  \end{tikzpicture}
\end{textblock*}
\begin{textblock*}{63.59mm}(8mm,288.27mm)
  {\sffamily\bfseries\color{accent}\fontsize{13}{16}\selectfont SAT 1200+}
\end{textblock*}
\begin{textblock*}{10mm}(188mm,288.27mm)
  \centering{\sffamily\color{accent}\fontsize{14}{16}\selectfont #1}
\end{textblock*}
}

\arrayrulecolor{accent}
\setlength{\arrayrulewidth}{0.7pt}

\begin{document}

\topicheader{P}{SAT iboralari}
\pagekicker{Phrases -- Tez-tez uchraydigan iboralar}
\pagefooter{85}
\mbox{}%
\begin{textblock*}{178mm}(16mm,33mm)
{\fontsize{10}{13}\selectfont\color{ink}
Quyidagi iboralar College Board'ning rasmiy savollar bazasida eng ko'p uchraydigan, mavzudan qat'i nazar takrorlanadigan tayyor so'z birikmalaridir. Chiziqcha (\textit{...}) o'rniga aniq son, o'zgaruvchi yoki ifoda qo'yiladi. Chastotasi bo'yicha tartiblangan -- eng ko'p uchraydigani birinchi.
}
\end{textblock*}
\begin{textblock*}{178mm}(16mm,46mm)
{\fontsize{11}{15}\selectfont
\begin{tabular}{>{\centering\arraybackslash}p{8mm}|p{89mm}|p{67mm}}
\hline
\rowcolor{accent}
{\color{white}\sffamily\bfseries \#} & {\color{white}\sffamily\bfseries English} & {\color{white}\sffamily\bfseries O'zbekcha} \\ \hline
1 & the value of & \textit{...ning qiymati} \\ \hline
\rowcolor{badge!18}
2 & given that & \textit{...berilgan bo'lsa} \\ \hline
3 & the number of & \textit{...soni} \\ \hline
\rowcolor{badge!18}
4 & both sides of & \textit{...ning ikkala tomoni} \\ \hline
5 & it follows that & \textit{bundan kelib chiqadiki} \\ \hline
\rowcolor{badge!18}
6 & is equivalent to & \textit{...ga teng (ekvivalent)} \\ \hline
7 & the graph of & \textit{...ning grafigi} \\ \hline
\rowcolor{badge!18}
8 & the length of & \textit{...ning uzunligi} \\ \hline
9 & represents the & \textit{...ni ifodalaydi} \\ \hline
\rowcolor{badge!18}
10 & in the xy-plane & \textit{xy-tekislikda} \\ \hline
11 & the correct answer is & \textit{to'g'ri javob -- bu} \\ \hline
\rowcolor{badge!18}
12 & each side of & \textit{har bir tomoni} \\ \hline
13 & which of the following & \textit{quyidagilardan qaysi biri} \\ \hline
\rowcolor{badge!18}
14 & dividing both sides & \textit{ikkala tomonni bo'lish} \\ \hline
15 & the measure of & \textit{...ning o'lchami} \\ \hline
\rowcolor{badge!18}
16 & subtracting from both sides & \textit{ikkala tomondan ayirish} \\ \hline
17 & the y-intercept & \textit{Y o'qidagi kesilgan nuqta} \\ \hline
\rowcolor{badge!18}
18 & the sum of & \textit{...yig'indisi} \\ \hline
19 & greater than & \textit{...dan katta} \\ \hline
\rowcolor{badge!18}
20 & less than & \textit{...dan kichik} \\ \hline
21 & rather than & \textit{...o'rniga} \\ \hline
\rowcolor{badge!18}
22 & the total number of & \textit{jami soni} \\ \hline
23 & there are & \textit{mavjud} \\ \hline
\rowcolor{badge!18}
24 & can be represented by & \textit{...bilan ifodalanishi mumkin} \\ \hline
25 & in the form & \textit{...ko'rinishida} \\ \hline
\rowcolor{badge!18}
26 & per hour & \textit{soatiga} \\ \hline
27 & at least & \textit{kamida} \\ \hline
\rowcolor{badge!18}
28 & adding to both sides & \textit{ikkala tomonga qo'shish} \\ \hline
29 & more than & \textit{...dan ortiq} \\ \hline
\rowcolor{badge!18}
30 & is approximately & \textit{taxminan ...ga teng} \\ \hline
\end{tabular}
}
\end{textblock*}
\newpage
\topicheader{P}{SAT iboralari}
\pagekicker{Phrases -- Tez-tez uchraydigan iboralar}
\pagefooter{86}
\mbox{}%
\begin{textblock*}{178mm}(16mm,33mm)
{\fontsize{11}{15}\selectfont
\begin{tabular}{>{\centering\arraybackslash}p{8mm}|p{89mm}|p{67mm}}
\hline
\rowcolor{accent}
{\color{white}\sffamily\bfseries \#} & {\color{white}\sffamily\bfseries English} & {\color{white}\sffamily\bfseries O'zbekcha} \\ \hline
31 & passes through the point & \textit{...nuqtadan o'tadi} \\ \hline
\rowcolor{badge!18}
32 & the slope of the line & \textit{to'g'ri chiziqning qiyaligi} \\ \hline
33 & corresponds to & \textit{...ga mos keladi} \\ \hline
\rowcolor{badge!18}
34 & in terms of & \textit{...orqali (ifodalangan)} \\ \hline
35 & a solution to & \textit{...ning yechimi} \\ \hline
\rowcolor{badge!18}
36 & the ratio of & \textit{...nisbati} \\ \hline
37 & according to the & \textit{...ga ko'ra} \\ \hline
\rowcolor{badge!18}
38 & divided by & \textit{...ga bo'lingan} \\ \hline
39 & can be written as & \textit{...shaklida yozilishi mumkin} \\ \hline
\rowcolor{badge!18}
40 & increased by & \textit{...ga oshirilgan} \\ \hline
41 & is represented by & \textit{...bilan ifodalanadi} \\ \hline
\rowcolor{badge!18}
42 & is closest to & \textit{...ga eng yaqin} \\ \hline
43 & intersects the & \textit{...ni kesib o'tadi} \\ \hline
\rowcolor{badge!18}
44 & results in & \textit{natijasida ...kelib chiqadi} \\ \hline
45 & based on the & \textit{...asosida} \\ \hline
\rowcolor{badge!18}
46 & is given by & \textit{...orqali beriladi} \\ \hline
47 & the sum of the measures & \textit{o'lchamlar yig'indisi} \\ \hline
\rowcolor{badge!18}
48 & for which & \textit{qaysi ...uchun} \\ \hline
49 & consists of & \textit{...dan iborat} \\ \hline
\rowcolor{badge!18}
50 & instead of & \textit{...o'rniga} \\ \hline
51 & the x-intercept & \textit{X o'qidagi kesilgan nuqta} \\ \hline
\rowcolor{badge!18}
52 & the ratio of the & \textit{...ning nisbati} \\ \hline
53 & the product of & \textit{...ko'paytmasi} \\ \hline
\rowcolor{badge!18}
54 & is parallel to & \textit{...ga parallel} \\ \hline
55 & has no solution & \textit{yechimga ega emas} \\ \hline
\rowcolor{badge!18}
56 & to the nearest & \textit{...gacha yaxlitlab} \\ \hline
57 & at most & \textit{ko'pi bilan} \\ \hline
\rowcolor{badge!18}
58 & the difference between & \textit{...orasidagi farq} \\ \hline
59 & there is & \textit{mavjud} \\ \hline
\rowcolor{badge!18}
60 & multiplied by & \textit{...ga ko'paytirilgan} \\ \hline
\end{tabular}
}
\end{textblock*}
\newpage
\topicheader{P}{SAT iboralari}
\pagekicker{Phrases -- Tez-tez uchraydigan iboralar}
\pagefooter{87}
\mbox{}%
\begin{textblock*}{178mm}(16mm,33mm)
{\fontsize{11}{15}\selectfont
\begin{tabular}{>{\centering\arraybackslash}p{8mm}|p{89mm}|p{67mm}}
\hline
\rowcolor{accent}
{\color{white}\sffamily\bfseries \#} & {\color{white}\sffamily\bfseries English} & {\color{white}\sffamily\bfseries O'zbekcha} \\ \hline
61 & corresponding angles & \textit{mos burchaklar} \\ \hline
\rowcolor{badge!18}
62 & such that & \textit{shundayki} \\ \hline
63 & the range of & \textit{...ning qiymatlar sohasi} \\ \hline
\rowcolor{badge!18}
64 & has infinitely many solutions & \textit{cheksiz ko'p yechimga ega} \\ \hline
65 & less than or equal to & \textit{...dan kichik yoki teng} \\ \hline
\rowcolor{badge!18}
66 & is perpendicular to & \textit{...ga perpendikulyar} \\ \hline
67 & per day & \textit{kuniga} \\ \hline
\rowcolor{badge!18}
68 & must be true & \textit{albatta to'g'ri bo'lishi kerak} \\ \hline
69 & no more than & \textit{ko'pi bilan / ...dan oshmaydi} \\ \hline
\rowcolor{badge!18}
70 & data shown & \textit{keltirilgan ma'lumotlar} \\ \hline
71 & vertical angles & \textit{vertikal burchaklar} \\ \hline
\rowcolor{badge!18}
72 & if and only if & \textit{faqat va faqat ...bo'lgandagina} \\ \hline
73 & which of the following could be & \textit{quyidagilardan qaysi biri ...bo'lishi mumkin} \\ \hline
\rowcolor{badge!18}
74 & is defined as & \textit{...deb ataladi} \\ \hline
75 & shown in the table & \textit{jadvalda ko'rsatilgan} \\ \hline
\rowcolor{badge!18}
76 & at a constant rate & \textit{doimiy sur'atda} \\ \hline
77 & the quadratic formula & \textit{kvadrat tenglama ildizlari formulasi} \\ \hline
\rowcolor{badge!18}
78 & has exactly one solution & \textit{aynan bitta yechimga ega} \\ \hline
79 & per year & \textit{yiliga} \\ \hline
\rowcolor{badge!18}
80 & not equal to & \textit{...ga teng emas} \\ \hline
81 & greater than or equal to & \textit{...dan katta yoki teng} \\ \hline
\rowcolor{badge!18}
82 & decreased by & \textit{...ga kamaytirilgan} \\ \hline
83 & a system of equations & \textit{tenglamalar sistemasi} \\ \hline
\rowcolor{badge!18}
84 & in order to & \textit{...uchun} \\ \hline
85 & to the nearest whole number & \textit{eng yaqin butun songacha} \\ \hline
\rowcolor{badge!18}
86 & is a factor of & \textit{...ning ko'paytuvchisi} \\ \hline
87 & solve for & \textit{...ni toping} \\ \hline
\rowcolor{badge!18}
88 & on average & \textit{o'rtacha hisobda} \\ \hline
89 & for all values & \textit{barcha qiymatlar uchun} \\ \hline
\rowcolor{badge!18}
90 & as shown & \textit{ko'rsatilganidek} \\ \hline
\end{tabular}
}
\end{textblock*}
\newpage
\topicheader{P}{SAT iboralari}
\pagekicker{Phrases -- Tez-tez uchraydigan iboralar}
\pagefooter{88}
\mbox{}%
\begin{textblock*}{178mm}(16mm,33mm)
{\fontsize{11}{15}\selectfont
\begin{tabular}{>{\centering\arraybackslash}p{8mm}|p{89mm}|p{67mm}}
\hline
\rowcolor{accent}
{\color{white}\sffamily\bfseries \#} & {\color{white}\sffamily\bfseries English} & {\color{white}\sffamily\bfseries O'zbekcha} \\ \hline
91 & alternate interior angles & \textit{almashinuvchi ichki burchaklar} \\ \hline
\rowcolor{badge!18}
92 & satisfies the equation & \textit{tenglamani qanoatlantiradi} \\ \hline
93 & is decreasing & \textit{kamayib bormoqda (kamayuvchi funksiya)} \\ \hline
\rowcolor{badge!18}
94 & rounded to the nearest & \textit{...gacha yaxlitlangan} \\ \hline
95 & the difference of & \textit{...ayirmasi} \\ \hline
\rowcolor{badge!18}
96 & complete the square & \textit{to'liq kvadratga keltirish} \\ \hline
97 & is a linear function & \textit{chiziqli funksiya} \\ \hline
\rowcolor{badge!18}
98 & as well as & \textit{shuningdek} \\ \hline
99 & is one of the & \textit{...lardan biri} \\ \hline
\rowcolor{badge!18}
100 & at a rate of & \textit{...sur'atda} \\ \hline
101 & is increasing & \textit{ortib bormoqda (o'suvchi funksiya)} \\ \hline
\rowcolor{badge!18}
102 & is a quadratic function & \textit{kvadrat funksiya} \\ \hline
103 & could be the value & \textit{qiymati bo'lishi mumkin} \\ \hline
\rowcolor{badge!18}
104 & shown in the graph & \textit{grafikda ko'rsatilgan} \\ \hline
105 & over time & \textit{vaqt o'tishi bilan} \\ \hline
\rowcolor{badge!18}
106 & is exactly & \textit{aynan ...ga teng} \\ \hline
107 & as a function of & \textit{...ning funksiyasi sifatida} \\ \hline
\rowcolor{badge!18}
108 & reaches a maximum & \textit{maksimumga erishadi} \\ \hline
109 & changes by & \textit{...ga o'zgaradi} \\ \hline
\rowcolor{badge!18}
110 & is a solution of & \textit{...ning yechimi hisoblanadi} \\ \hline
111 & with respect to & \textit{...ga nisbatan} \\ \hline
\rowcolor{badge!18}
112 & the discriminant of & \textit{...ning diskriminanti} \\ \hline
113 & in addition to & \textit{...dan tashqari} \\ \hline
\rowcolor{badge!18}
114 & lies on the line & \textit{chiziq ustida yotadi} \\ \hline
115 & is a multiple of & \textit{...ning karralisi} \\ \hline
\rowcolor{badge!18}
116 & distribute the & \textit{...ni taqsimlang (distributivlik qoidasini qo'llang)} \\ \hline
117 & depends on & \textit{...ga bog'liq} \\ \hline
\rowcolor{badge!18}
118 & combine like terms & \textit{o'xshash hadlarni birlashtiring} \\ \hline
119 & which of the following best describes & \textit{quyidagilardan qaysi biri eng yaxshi ifodalaydi} \\ \hline
\rowcolor{badge!18}
120 & regardless of & \textit{...dan qat'i nazar} \\ \hline
\end{tabular}
}
\end{textblock*}
\newpage
\topicheader{P}{SAT iboralari}
\pagekicker{Phrases -- Tez-tez uchraydigan iboralar}
\pagefooter{89}
\mbox{}%
\begin{textblock*}{178mm}(16mm,33mm)
{\fontsize{11}{15}\selectfont
\begin{tabular}{>{\centering\arraybackslash}p{8mm}|p{89mm}|p{67mm}}
\hline
\rowcolor{accent}
{\color{white}\sffamily\bfseries \#} & {\color{white}\sffamily\bfseries English} & {\color{white}\sffamily\bfseries O'zbekcha} \\ \hline
121 & is tangent to & \textit{...ga urinadi} \\ \hline
\rowcolor{badge!18}
122 & is an exponential function & \textit{ko'rsatkichli funksiya} \\ \hline
123 & is undefined & \textit{aniqlanmagan} \\ \hline
\rowcolor{badge!18}
124 & in a given & \textit{berilgan ...da} \\ \hline
125 & for some value & \textit{qandaydir qiymat uchun} \\ \hline
\rowcolor{badge!18}
126 & expressed in terms of & \textit{...orqali ifodalangan} \\ \hline
127 & due to & \textit{...sababli} \\ \hline
\rowcolor{badge!18}
128 & made up of & \textit{...dan tashkil topgan} \\ \hline
129 & factor the expression & \textit{ifodani ko'paytuvchilarga ajrating} \\ \hline
\rowcolor{badge!18}
130 & written in terms of & \textit{...orqali yozilgan} \\ \hline
131 & the domain of & \textit{...ning aniqlanish sohasi} \\ \hline
\rowcolor{badge!18}
132 & substitute the value & \textit{qiymatni o'rniga qo'ying} \\ \hline
133 & simplify the expression & \textit{ifodani soddalashtiring} \\ \hline
\rowcolor{badge!18}
134 & remains constant & \textit{o'zgarmas qoladi} \\ \hline
135 & real number solutions & \textit{haqiqiy sonli yechimlar} \\ \hline
\rowcolor{badge!18}
136 & reaches a minimum & \textit{minimumga erishadi} \\ \hline
137 & on the interval & \textit{oraliqda} \\ \hline
\rowcolor{badge!18}
138 & none of the above & \textit{yuqoridagilarning hech biri} \\ \hline
139 & is proportional to & \textit{...ga proporsional} \\ \hline
\rowcolor{badge!18}
140 & inversely proportional to & \textit{...ga teskari proporsional} \\ \hline
141 & in common & \textit{umumiy} \\ \hline
\rowcolor{badge!18}
142 & exists such that & \textit{shunday ...mavjudki} \\ \hline
143 & evaluate the expression & \textit{ifodaning qiymatini hisoblang} \\ \hline
\rowcolor{badge!18}
144 & directly proportional to & \textit{...ga to'g'ri proporsional} \\ \hline
145 & complex solutions & \textit{kompleks yechimlar} \\ \hline
\rowcolor{badge!18}
146 & as represented & \textit{ifodalanganidek} \\ \hline
\end{tabular}
}
\end{textblock*}
\newpage
\end{document}
